\documentclass[12pt,a4paper]{article}

\usepackage[utf8]{inputenc}
\usepackage[portuguese]{babel}
\usepackage{amsmath}
\usepackage{graphicx}
\usepackage{booktabs}
\usepackage{hyperref}
\usepackage{geometry}
\geometry{a4paper, margin=1in}

\title{Relatório Final: Avaliação do Modelo TabICL v2}
\author{Disciplina: Aprendizagem de Máquina}
\date{\today}

\begin{document}

\maketitle

\begin{abstract}
Este projeto teve como objetivo avaliar o desempenho do foundation model \textbf{TabICL v2}, focado em In-Context Learning (ICL) para dados tabulares. O modelo foi testado em 30 datasets extraídos do benchmark TabArena-v0.1 e comparado com métodos estado-da-arte baseados em árvores (LightGBM, XGBoost e CatBoost). Os resultados indicam que o TabICL v2 é altamente superior em capacidade preditiva (AUC de 0.928 vs 0.907 do melhor baseline). No entanto, sua arquitetura baseada em atenção possui limitações de escalabilidade para datasets com alta dimensionalidade e grande volume de amostras, resultando em falhas de memória (Out-Of-Memory) em 4 dos 30 datasets testados.
\end{abstract}

\section{Metodologia}

A avaliação seguiu um rigoroso pipeline de testes padronizado:
\begin{itemize}
    \item \textbf{Datasets:} Foram selecionados 30 datasets multiclasse e binários do TabArena-v0.1 via OpenML.
    \item \textbf{Divisão (Split):} Utilizou-se um particionamento estratificado de 70\% para treino (usado como contexto para o modelo ICL) e 30\% para teste, garantindo a reprodutibilidade com a semente (seed) fixa em 42.
    \item \textbf{Pré-processamento:} O pipeline foi aprimorado para codificar classes e features categóricas de forma automatizada. Também implementou-se a imputação de valores nulos (NaN) usando a mediana, garantindo que os modelos baseline tradicionais não falhassem durante o treinamento. O TabICL é capaz de processar essas anomalias de forma nativa.
    \item \textbf{Métricas:} Foram extraídas as métricas AUC OvO (métrica primária), Acurácia, G-Mean e Cross-Entropy.
\end{itemize}

\section{Resultados Quantitativos}

\subsection{Médias Agregadas}
Nas avaliações agregadas (excluindo os 4 datasets com limitações de hardware discutidos na Seção \ref{sec:limitacoes}), o TabICL v2 obteve a melhor performance preditiva de forma consistente em todas as métricas analisadas.

\begin{table}[h]
\centering
\caption{Desempenho médio dos modelos testados}
\vspace{0.2cm}
\begin{tabular}{@{}lccc@{}}
\toprule
\textbf{Modelo} & \textbf{AUC OvO} & \textbf{Acurácia} & \textbf{G-Mean} \\ \midrule
\textbf{TabICL v2} & \textbf{0.9287} & \textbf{0.8787} & \textbf{0.8003} \\
CatBoost & 0.9072 & 0.8524 & 0.7312 \\
XGBoost & 0.9016 & 0.8530 & 0.7519 \\
LightGBM & 0.8995 & 0.8525 & 0.6909 \\ \bottomrule
\end{tabular}
\end{table}

\subsection{Análise Estatística (Friedman e Nemenyi)}
O teste de Friedman confirmou uma diferença estatisticamente significativa entre os classificadores. O ranking médio (Mean Rank) sobre os múltiplos datasets isola a performance do TabICL como amplamente superior:

\begin{enumerate}
    \item \textbf{TabICL v2:} 1.17
    \item \textbf{CatBoost:} 2.65
    \item \textbf{LightGBM:} 2.94
    \item \textbf{XGBoost:} 3.23
\end{enumerate}

O diagrama de diferença crítica (CD diagram) gerado a partir do Nemenyi post-hoc confirma que o modelo TabICL se separa de forma estatisticamente significativa do grupo de baselines baseados em árvores de decisão.

\subsection{Teste Bayesiano com ROPE}
O teste bayesiano com uma Região de Equivalência Prática (ROPE) configurada para uma margem de 0.01 na métrica AUC comparou diretamente o TabICL contra o LightGBM, concluindo:
\begin{itemize}
    \item \textbf{60.0\%} de probabilidade estatística de o TabICL ser práticamente superior (margem de vitória $> 0.01$).
    \item \textbf{40.0\%} de probabilidade de ambos serem considerados estatisticamente equivalentes.
    \item \textbf{0.0\%} de probabilidade de o TabICL ser inferior ao baseline.
\end{itemize}

\section{Observações e Limitações: Falhas por OOM}
\label{sec:limitacoes}

Apesar do desempenho excepcional em métricas de classificação teóricas, a arquitetura foundation do TabICL v2 revelou um sério gargalo computacional durante o experimento prático. 

Em \textbf{4 dos 30 datasets avaliados} (\textit{isolet}, \textit{jm1}, \textit{adult}, e \textit{Bioresponse}), o modelo encontrou falhas críticas de hardware e teve sua execução abortada pelo sistema operacional. O cálculo de inferência em In-Context Learning aplica uma matriz de atenção entre todas as amostras de treino (contexto) e as amostras de teste. Isso resulta em uma complexidade assintótica de tempo e espaço de $\mathcal{O}(n^2 \cdot p)$. Nesses 4 datasets, que combinam altíssima dimensionalidade de features com milhares de instâncias, a exigência de memória RAM/VRAM escalou para dezenas de gigabytes, causando travamentos por falta de memória (Out-Of-Memory).

\section{Conclusão}

O TabICL v2 representa um avanço significativo em precisão out-of-the-box para problemas tabulares, superando até mesmo os modelos estado-da-arte baseados em Gradient Boosting sem a necessidade explícita de sintonia fina (hyperparameter tuning). 

É o modelo ideal para datasets pequenos e médios ($< 10.000$ instâncias) onde o ganho de AUC compensa largamente o tempo computacional extra da inferência. No entanto, para enormes volumes de dados de produção, modelos como CatBoost e LightGBM continuam sendo alternativas substancialmente superiores do ponto de vista de infraestrutura e engenharia de software devido à sua extrema eficiência e estabilidade.

\end{document}
